%%
% 第5章 智能旅行规划系统的设计与实现
%%

\chapter{智能旅行路线推荐系统的设计与实现}
\label{chap:system}

%要不要把用户意图可视化作为一个卖点，还是说路线可视化作为一个核心点？
第三章提出了基于知识蒸馏与显式意图增强的单城市旅行路线推荐框架EKD-Trip，该框架将历史轨迹中的隐式偏好蒸馏至查询编码器，并引入四种旅行模式（接近、远离、U型、不规则）的分类任务以增强用户空间移动意图的建模。第四章进一步提出了融合多维用户偏好意图的跨城市旅行路线推荐模型CrossTrip，该模型利用大语言模型提取高层语义偏好，经多头解耦机制分解用户偏好，并结合自适应迁移门控与城市群体记忆实现跨城市冷启动场景下的个性化推荐。本章在上述两章工作的基础上，将旅行路线推荐算法落地为一个面向实际旅行规划场景的智能系统TripManner。TripManner提供自然语言交互式路线规划、用户偏好意图可视化和路线地图展示等功能，用户可在对话中获取个性化的路线推荐，并直观理解推荐背后的意图推断逻辑，从而建立对推荐结果的信任。
%利用TripManner系统，减轻用户在出行时做攻略的时间和精力成本，同时通过路线同步可视化以及随时调整机制，让

\section{TripManner系统需求分析}
\label{sec:ch5_requirement}

随着旅游产业的快速发展和移动互联网的普及，越来越多的用户依赖在线服务来规划旅行路线。然而，现有的旅行规划应用往往存在推荐路线单一、无法感知用户深层偏好、交互方式不够智能等问题。本文在EKD-Trip和CrossTrip两项工作的基础上，设计并实现了智能旅行路线推荐系统TripManner（Intelligent Travel Route Recommendation System）。该系统利用历史签到数据挖掘用户出行意图，并以可视化界面辅助用户理解推荐逻辑。为将论文算法落地为可信、可解释的旅行规划服务，系统应至少满足以下需求：

（1）系统需支持基于自然语言的交互式旅行规划，用户能够通过对话方式描述旅行需求，系统自动解析用户意图并生成个性化路线推荐。

（2）系统需集成本文第三章和第四章提出的路线推荐算法，当用户查询涉及算法数据集覆盖的城市时，能够调用对应算法进行意图感知的路线推荐；对于其他城市，系统需提供基于大语言模型的通用路线规划能力。

（3）系统需提供用户偏好意图的可视化分析功能，将算法推理过程中识别的旅行模式、偏好因子、迁移权重等信息以直观的图表方式呈现，增强系统的可解释性和用户信任度。

（4）系统界面应支持流式输出和动态路线渲染，路线规划结果应以打字机效果逐步呈现，同时在地图上动态绘制路线，提升用户的交互体验。

（5）系统界面设计应简洁直观，具备良好的可视化展示能力，使用户能够快速理解推荐结果及其背后的意图推断逻辑。


\subsection{功能性需求分析}
\label{sec:ch5_func_req}

根据上述系统需求分析，将TripManner系统的功能性需求划分为六个模块：用户管理模块、智能路线规划模块、意图可视化模块、地图可视化模块、对话交互模块和数据管理模块。系统用例图如图\ref{fig:ch5_usecase}所示，下面将具体介绍各模块负责的功能。

\begin{figure}[hbt]
  \centering
  % TODO: 插入系统用例图
  % 内容：展示用户和管理员两个角色与六大功能模块的交互关系
  \fbox{\parbox{0.8\textwidth}{\centering\vspace{3cm}系统用例图占位\vspace{3cm}}}
  \caption{TripManner系统用例图}
  \label{fig:ch5_usecase}
\end{figure}

（1）用户管理模块。该模块处理用户认证与个人信息管理，包括：用户注册与登录（用户名须唯一），支持修改密码、个人信息和权限；用户偏好设置，允许配置旅行兴趣类别、出行风格、同行人类型和预算级别等参数；历史规划记录的查看与管理。

（2）智能路线规划模块。该模块是系统的核心，接收用户旅行查询请求后由路由决策引擎自动选择算法进行路线规划。系统支持三种规划路径：目的城市属于单城市数据集（Glasgow、Osaka、Toronto、Tokyo）时调用EKD-Trip算法；查询涉及跨城市场景且目标城市属于跨城市数据集（New York、Los Angeles、San Francisco）时调用CrossTrip算法；其他城市则由DeepSeek大语言模型Agent生成路线。

（3）意图可视化模块。该模块将算法推理过程中识别的用户偏好以可视化方式呈现。单城市场景下展示EKD-Trip识别的旅行模式分类结果及用户到目的地的距离变化曲线；跨城市场景下展示CrossTrip解耦的四维偏好因子雷达图、偏好迁移权重、城市群体偏好分布以及个人偏好与群体偏好的融合比例。

（4）地图可视化模块。该模块基于Leaflet地图引擎与OpenStreetMap实现交互式地图展示，支持路线连线绘制、POI标记与详情弹窗、起终点高亮标注及自动视窗适配等功能。

（5）对话交互模块。该模块实现基于自然语言的对话式旅行规划交互。当用户输入的查询信息不完整时，系统通过引导式对话逐步补全查询五元组（出发城市、出发时间、目的城市、结束时间、途经点数）。路线规划结果经Server-Sent Events（SSE）流式推送至前端，以打字机效果逐步渲染文本，同时在地图上同步绘制路线。

（6）数据管理模块。该模块负责系统数据的存储与管理，包括数据集城市的POI信息管理、用户历史聊天记录存储、用户签到历史数据存储以及系统配置参数管理。

\subsection{非功能性需求分析}
\label{sec:ch5_nonfunc_req}

TripManner系统除需满足上述功能性需求外，还应具备以下非功能性需求，以保证系统在实际使用中的可靠性和用户体验：

（1）响应时间：系统操作响应时间应小于500毫秒，页面加载时间应小于3秒，大语言模型Agent模式下首个SSE事件的响应时间应小于2秒，处理过程中以加载动画提示用户。

（2）流式交互体验：文本流式输出的逐Token延迟应小于100毫秒，POI标记在地图上的动态添加间隔应控制在300至800毫秒之间，以模拟自然的思考节奏，提升交互流畅感。

（3）可用性：系统需要保证全时可用，所有功能模块均应可正常访问和操作，确保其年度可用性不低于99%。

（4）可维护性：系统代码结构应清晰，遵循前后端分离的模块化设计原则，各功能模块之间保持低耦合，便于后期的维护与功能扩展。

（5）可扩展性：系统应支持灵活扩展，能适应未来需求接入。新增数据源与接口时应平稳过渡，减少对现有功能的影响。

（6）安全性：系统应实施严格的身份验证、身份管理和权限控制措施，确保用户数据的机密性、完整性和可用性。用户密码采用BCrypt加密存储，大语言模型API密钥不在前端代码中暴露。

（7）兼容性：系统前端应兼容Chrome、Firefox和Edge等主流浏览器，流式输出功能依赖的EventSource API需确保跨浏览器兼容性。

（8）模型关键性指标：为保障推荐路线的实际效果，本文提出量化指标约束推荐准确性。推荐地点精确度方面，$F_1$指标不低于50\%，即用户实际访问的兴趣点至少有50\%出现在推荐结果中。地理空间范围方面，模型平均误差控制在半径5km以内，即实际路线与预测路线之间的空间偏差不超过5km。

\section{架构设计}
\label{sec:ch5_architecture}

本系统采用分层架构设计思想，自上而下划分为用户层、平台层和数据层，如图\ref{fig:ch5_arch}所示。

\begin{figure}[hbt]
  \centering
  % TODO: 插入系统总体架构图
  % 内容：三层架构图——用户层（六大功能模块）、平台层（路由决策引擎+三条算法路径）、数据层（MySQL各表）
  \fbox{\parbox{0.85\textwidth}{\centering\vspace{4cm}系统总体架构图占位\vspace{4cm}}}
  \caption{TripManner系统总体架构图}
  \label{fig:ch5_arch}
\end{figure}

用户层承载系统的六大功能模块界面，包括对话交互、地图可视化、意图可视化面板、用户管理、历史记录和系统管理。用户通过自然语言对话提交旅行查询，系统以流式文本、动态地图路线和意图分析图表的形式返回规划结果。

平台层是系统的核心计算层，其核心组件为路由决策引擎。该引擎根据用户查询的目的地城市判断应调用的算法路径：目的城市属于单城市数据集时调用EKD-Trip算法服务；涉及跨城市场景且目标城市在跨城市数据集范围内时调用CrossTrip算法服务；否则调用DeepSeek大语言模型Agent进行通用路线规划。

数据层负责系统数据的持久化存储，基于MySQL关系型数据库管理用户数据、POI信息、路线规划记录、对话历史等结构化数据。

为满足系统的高效性与可维护性需求，本系统采用前后端分离架构，如图\ref{fig:ch5_techstack}所示，前端与后端通过RESTful API和SSE流式接口进行通信。

\begin{figure}[hbt]
  \centering
  % TODO: 插入前后端技术框架图
  % 内容：左侧前端技术栈（Vue3/Element Plus/Leaflet/ECharts/Axios）、右侧后端技术栈（FastAPI/SQLAlchemy/BCrypt）、
  % 中间算法支撑（EKD-Trip/CrossCityLLMCPR/DeepSeek API）、底部数据存储（MySQL）
  \fbox{\parbox{0.85\textwidth}{\centering\vspace{3.5cm}前后端技术框架图占位\vspace{3.5cm}}}
  \caption{TripManner前后端技术框架}
  \label{fig:ch5_techstack}
\end{figure}

前端基于Vue 3框架，使用Vite构建，结合Element Plus进行UI组件开发。地图可视化采用Leaflet引擎配合OpenStreetMap瓦片服务，无需额外的地图授权密钥。意图可视化图表（雷达图、折线图、饼图等）由ECharts库渲染。前后端数据通信使用Axios HTTP客户端，流式输出通过浏览器原生的EventSource API接收SSE事件流。

后端采用Python FastAPI框架，原生支持异步请求处理和SSE流式响应。数据库操作通过SQLAlchemy ORM实现，用户认证采用密码加密校验机制。后端封装了三个算法服务接口：EKD-Trip推理接口、CrossTrip推理接口和DeepSeek API调用接口。

算法支撑层集成了本文第三章和第四章提出的旅行路线推荐算法，以及DeepSeek大语言模型的API调用能力。EKD-Trip算法负责单城市场景下基于知识蒸馏和旅行模式预测的路线推荐；CrossTrip算法负责跨城市场景下融合多维偏好意图的路线推荐；DeepSeek Agent负责数据集未覆盖城市的通用路线生成。

数据存储层采用MySQL 8.0关系型数据库，存储用户信息、POI数据、路线规划记录、对话历史等全部结构化数据。

\section{详细设计}
\label{sec:ch5_detail}

\subsection{用户管理模块}
\label{sec:ch5_user}

用户管理模块处理用户认证、个人信息管理和偏好设置，保障系统安全性和个性化服务能力。模块包含三个核心类：\texttt{UserManager}类、\texttt{SessionManager}类和\texttt{Authorization}类，如图\ref{fig:ch5_user_uml}所示。

\begin{figure}[hbt]
  \centering
  % TODO: 插入用户管理模块UML类图
  % 内容：UserManager（addUser/updateUser/deleteUser/queryUser）、SessionManager（storeSession/retrieveSession）、
  % Authorization（verifyCredentials/checkPermission），三者之间的关联关系
  \fbox{\parbox{0.7\textwidth}{\centering\vspace{3cm}用户管理模块UML类图占位\vspace{3cm}}}
  \caption{用户管理模块UML设计}
  \label{fig:ch5_user_uml}
\end{figure}

\texttt{UserManager}类管理用户信息，包括创建、更新、删除及查询操作。注册时密码经BCrypt算法加密后存入数据库，明文密码不被保留。用户偏好以JSON格式存储，包含旅行兴趣类别（\texttt{interested\_categories}，如文化古迹、自然风光、美食等）、出行风格（\texttt{travel\_style}，如深度游、打卡游）、同行人类型（\texttt{companion\_type}，如独自、家庭、朋友）以及预算级别（\texttt{budget\_level}）。这些偏好在路线规划时作为用户画像输入到算法或大语言模型中。

\texttt{Authorization}类负责基于数据库的用户身份验证与权限校验。系统采用直接校验机制：在用户登录或发起敏感操作时，直接查询数据库比对账号凭证（如用户名与加密密码）的有效性。该机制通过实时查库确保身份数据的强一致性，在简化系统架构的同时满足了基础场景下的安全认证需求。

\texttt{SessionManager}类负责用户会话的管理，记录用户的登录状态和活跃会话信息，确保每次请求都能进行有效的身份认证。

\subsection{智能路线规划模块}
\label{sec:ch5_route}

智能路线规划模块是TripManner系统的核心，其设计重点在于路由决策引擎。引擎根据用户查询的目的地城市选择最合适的算法路径，流程如图\ref{fig:ch5_route_decision}所示。

\begin{figure}[hbt]
  \centering
  % TODO: 插入路由决策引擎流程图
  % 内容：用户查询输入 → 提取目的城市 → 判断是否属于单城市数据集/跨城市数据集/其他
  % → 分别调用EKD-Trip/CrossCityLLMCPR/DeepSeek Agent → 统一输出格式
  % 中间穿插路由决策判断
  \fbox{\parbox{0.85\textwidth}{\centering\vspace{4cm}路由决策引擎流程图占位\vspace{4cm}}}
  \caption{路由决策引擎流程图}
  \label{fig:ch5_route_decision}
\end{figure}

路由决策引擎接收用户的旅行查询五元组$Q = (c_d, t_s, c_a, t_e, N)$作为输入，其中$c_d$为出发地点（隐含了城市信息），$t_s$为出发时间，$c_a$为目的地点（隐含了城市信息），$t_e$为结束时间，$N$为期望途经点数。引擎首先提取出目的城市$c_a$，然后按照以下规则进行路由判断：

（1）若$c_a \in \mathcal{S}_{single} = \{\text{Glasgow}, \text{Osaka}, \text{Toronto}, \text{Tokyo}\}$，即目的城市属于第三章EKD-Trip算法的数据集覆盖城市，则调用EKD-Trip算法接口。该接口接收起点POI标识、终点POI标识、出发时间、结束时间和期望POI数量作为输入，返回推荐的POI序列及其详细信息，同时输出旅行模式分类结果（接近模式、远离模式、U型模式或不规则模式）和用户到目的地的距离变化序列。

（2）若查询涉及跨城市场景，且目标城市$c_a \in \mathcal{S}_{cross} = \{\text{New York}, \text{Los Angeles}, \text{San Francisco}\}$，即属于第四章CrossCityLLMCPR算法的数据集覆盖城市，则调用跨城市算法接口。该接口接收用户在源城市的签到历史、用户画像和查询条件作为输入，返回推荐的POI序列，同时输出$K=4$维偏好因子权重、偏好迁移系数$\alpha_k$、目标城市群体偏好向量、可靠性权重$\rho_g$以及个人偏好与群体偏好的融合门控$\eta$。

（3）若目的城市$c_a \notin \mathcal{S}_{single} \cup \mathcal{S}_{cross}$，即不在任何数据集的覆盖范围内，则调用DeepSeek大语言模型Agent。Agent接收用户的自然语言查询和偏好信息，通过设计好的提示词模板调用DeepSeek API生成结构化路线推荐，输出格式与算法接口一致，但不包含意图可视化数据。
%
%系统通过全局配置参数\texttt{USE\_MOCK\_DATA}控制算法调用的数据来源。当该参数为\texttt{True}时，EKD-Trip和CrossCityLLMCPR接口将直接从预设的Mock数据库中检索匹配的路线数据返回，而非实际调用算法推理服务；当参数为\texttt{False}时，系统将实际调用对应的算法推理接口。需要说明的是，DeepSeek Agent始终进行真实的API调用，不受Mock/Real开关影响，因为其不依赖本地部署的算法服务。

三条算法路径的输出均统一封装为标准化的路线结果格式，包含算法标识、城市信息、POI序列（含名称、类别、经纬度、建议停留时长）、意图数据（仅算法路径提供）以及生成置信度等字段，以便前端进行统一的展示和处理。

\subsection{意图可视化模块}
\label{sec:ch5_intent}

意图可视化模块是TripManner区别于一般旅行规划应用的特色所在，将算法推理过程中的用户偏好意图以图表形式直观展示，帮助用户理解推荐逻辑。模块根据路由决策引擎选择的算法路径渲染不同的可视化内容。

单城市场景下（EKD-Trip算法），意图可视化面板展示三部分内容：第一，旅行模式分类结果，以图标和文字标签展示系统识别出的旅行模式类型及其置信度。第三章定义的四种旅行模式各对应一种视觉符号：接近模式（Approaching）以绿色渐近箭头表示，说明用户倾向于沿途逐步靠近目的地；远离模式（Moving Away）以橙色远离箭头表示，说明用户先远离目的地再折返；U型模式（U-turn）以蓝色U型曲线表示，说明用户到目的地的距离先增后减；不规则模式（Irregular）以灰色散点表示，说明无明显空间趋势。第二，距离变化曲线图，使用ECharts折线图绘制，横轴为POI访问顺序，纵轴为当前POI到目的地的距离，直观反映旅行模式的空间含义。第三，模式说明文字，用简洁的自然语言向用户解释当前旅行模式的含义。

跨城市场景下（CrossTrip算法），意图可视化面板展示四部分内容：第一，四维偏好因子雷达图，使用ECharts雷达图绘制，四个轴分别对应$K=4$个正交偏好维度（如长期偏好、短期偏好、通勤偏好、娱乐偏好），展示各维度权重分布。第二，偏好迁移权重可视化，以带权重的箭头图或桑基图展示从源城市到目标城市四个偏好因子的迁移系数$\alpha_k$，连线粗细反映迁移权重大小，让用户看到哪些偏好被更多地迁移到了新城市。第三，城市群体偏好饼图，展示目标城市的群体偏好分布，配合可靠性得分$\rho_g$显示该城市数据的充分程度。第四，个人与群体偏好融合比例，以进度条形式展示融合门控$\eta$的值，让用户了解路线中个人偏好与群体偏好各占多大比重。

通用城市场景下（DeepSeek Agent），大语言模型不产出结构化的偏好意图数据，因此意图可视化面板不予显示，仅展示路线规划结果和Agent生成的文字说明。

\subsection{对话交互模块}
\label{sec:ch5_dialog}

对话交互模块实现基于自然语言的旅行规划交互，两大核心设计为引导式对话补全和流式输出与动态路线渲染。

\subsubsection{引导式对话补全}

实际使用中，用户的自然语言输入往往不包含路线规划所需的全部信息。例如用户可能只说"我想去东京玩"，而未给出出发时间、结束时间和期望途经点数等参数。为此系统设计了引导式对话补全机制，通过多轮对话逐步引导用户补全查询五元组。

补全流程如图\ref{fig:ch5_guide_flow}所示。系统将用户的自然语言输入和历史对话上下文发送至DeepSeek大语言模型进行语义解析，提取已有查询字段并识别缺失字段。若所有必填字段（目的地点、出发时间、结束时间）均已获取，五元组完整，直接进入路由决策引擎；若存在缺失字段，系统生成追问话术引导用户补充，如此循环直至五元组完整。

\begin{figure}[hbt]
  \centering
  % TODO: 插入引导式对话补全流程图
  % 内容：用户输入 → LLM解析提取字段 → 判断是否完整 → 完整则进入路由决策/不完整则生成追问 → 循环
  \fbox{\parbox{0.75\textwidth}{\centering\vspace{3.5cm}引导式对话补全流程图占位\vspace{3.5cm}}}
  \caption{引导式对话补全流程}
  \label{fig:ch5_guide_flow}
\end{figure}

五元组中各字段的引导策略如表\ref{tab:ch5_guide}所示。其中，目的地点、出发地点、出发和结束时间为必填字段，若用户未提及则系统首先追问并以示例引导用户回答；途经点数为可选字段，默认值为5个。

\begin{table}[hbt]
  \centering
  \caption{查询五元组字段引导策略}
  \label{tab:ch5_guide}
  \begin{tabular*}{0.92\textwidth}{@{\extracolsep{\fill}}cccc}
  \toprule
    字段 & 是否必填 & 默认值 & 引导话术示例 \\
  \midrule
  	出发地点 & 必填 & 无 & "您是从哪个地方出发的呢？" \\
    目的地点 & 必填 & 无 & "请问您的行程最后地点是哪里呢？" \\
    出发时间 & 必填 & 无 & "您计划什么时候出发？比如上午9点" \\
    结束时间 & 必填 & 无 & "大概玩到几点结束呢？" \\
    途经点数 & 可选 & 5 & "您希望途经几个景点？默认安排5个" \\
  \bottomrule
  \end{tabular*}
\end{table}

\subsubsection{流式输出与动态路线渲染}

路线规划结果采用Server-Sent Events（SSE）协议以流式方式推送至前端，实现逐步输出的打字机效果，同时地图上同步进行动态路线绘制。流式输出架构如图\ref{fig:ch5_sse}所示，后端FastAPI服务将路线结果封装为多个SSE事件逐步发送，前端通过浏览器原生的EventSource API监听事件流并实时渲染。

\begin{figure}[hbt]
  \centering
  % TODO: 插入SSE流式输出架构图
  % 内容：后端（算法/LLM调用 → 逐POI封装事件）→ SSE Stream → 前端（①对话区逐字渲染 ②地图区逐步添加POI ③意图区最后渲染图表）
  \fbox{\parbox{0.85\textwidth}{\centering\vspace{3cm}SSE流式输出架构图占位\vspace{3cm}}}
  \caption{SSE流式输出架构}
  \label{fig:ch5_sse}
\end{figure}

系统定义了七种SSE事件类型，如表\ref{tab:ch5_sse_events}所示。流式输出的时序过程为：首先发送\texttt{thinking}事件显示加载提示；然后逐个POI交替发送\texttt{route\_text}事件（文本描述）和\texttt{poi\_added}事件（POI坐标数据）；当所有POI输出完毕后发送\texttt{intent\_data}事件（仅算法路径）触发意图可视化面板的渲染；最后发送\texttt{done}事件结束流式输出。若在引导对话阶段检测到查询信息不完整，则发送\texttt{guide\_question}事件向用户追问。

\begin{table}[hbt]
  \centering
  \caption{SSE事件类型定义}
  \label{tab:ch5_sse_events}
  \begin{tabular*}{0.95\textwidth}{@{\extracolsep{\fill}}cccc}
  \toprule
    事件类型 & 触发时机 & 数据内容 & 前端行为 \\
  \midrule
    \texttt{thinking} & 开始处理 & 状态信息 & 显示加载提示 \\
    \texttt{guide\_question} & 需要追问 & 追问文本与缺失字段 & 对话区显示追问 \\
    \texttt{route\_text} & 路线文本片段 & 增量文本 & 逐字追加渲染 \\
    \texttt{poi\_added} & 新增一个POI & POI坐标与详情 & 地图添加标记和连线 \\
    \texttt{intent\_data} & 意图数据就绪 & 意图分析结果 & 渲染意图可视化面板 \\
    \texttt{done} & 输出完成 & 规划记录ID & 结束加载，保存记录 \\
    \texttt{error} & 出错 & 错误信息 & 显示错误提示 \\
  \bottomrule
  \end{tabular*}
\end{table}

前端动态渲染方面，POI标记出现时采用从小到大的弹入动画（300毫秒），路线连线从虚线到实线逐步绘制（500毫秒），意图可视化面板在所有POI输出完毕后淡入显示（400毫秒），地图视窗在每个新POI添加时平滑调整以包含已有POI（600毫秒）。这些动画使路线规划过程更加生动，提升了交互体验。

\subsection{地图可视化模块}
\label{sec:ch5_map}

地图可视化模块基于Leaflet开源地图引擎与OpenStreetMap瓦片服务，实现交互式地图展示。模块包含六项功能：第一，路线绘制，在地图上以带方向箭头的折线连接各推荐POI，展示访问路径和行进方向。第二，POI标记，每个推荐POI以带编号的标记显示在地图上，编号对应访问顺序。第三，POI详情弹窗，点击标记后弹出信息窗口，显示POI名称、类别、建议停留时长等信息。第四，起终点高亮，起点以绿色图标、终点以红色图标标注，便于快速识别路线起止位置。第五，自动视窗适配，路线生成后地图自动缩放至包含所有POI的最佳视窗。第六，地图交互操作，支持鼠标滚轮缩放和拖拽平移。

流式输出过程中，地图组件配合\texttt{poi\_added}事件动态更新：每收到一个新POI事件，地图即添加对应标记并绘制与上一POI之间的路线连线，视窗同步平滑调整以包含新增POI。这种逐步绘制的方式让用户随着文本描述的推进，看到路线在地图上逐步成形。

\subsection{数据管理模块}
\label{sec:ch5_data}

数据管理模块管理系统全部结构化数据，采用MySQL 8.0关系型数据库。数据库实体关系如图\ref{fig:ch5_er}所示，包含五张核心数据表。

\begin{figure}[hbt]
  \centering
  % TODO: 插入数据库ER图
  % 内容：users(1)──(N)routes, users(1)──(N)dialogs, users(1)──(N)user_checkins,
  % pois表独立, mock_routes表独立, user_checkins关联pois
  \fbox{\parbox{0.75\textwidth}{\centering\vspace{3cm}数据库ER图占位\vspace{3cm}}}
  \caption{TripManner数据库实体关系图}
  \label{fig:ch5_er}
\end{figure}

各数据表的设计如表\ref{tab:ch5_db}所示。\texttt{users}表存储用户基本信息和偏好设置，偏好数据以JSON格式存储以支持灵活扩展。\texttt{pois}表存储数据集城市的兴趣点信息，包含名称、类别、城市、经纬度坐标和热度等字段，其中城市字段建立索引以加速按城市查询。\texttt{routes}表记录每次路线规划的完整信息，包括原始查询输入、使用的算法标识、路线结果和意图数据，均以JSON格式存储；\texttt{algorithm\_used}字段为枚举类型，取值包括\texttt{EKD-Trip}、\texttt{CrossTrip}和\texttt{DeepSeek-Agent}。\texttt{dialogs}表存储对话记录，消息列表以JSON数组格式存储。\texttt{user\_checkins}表存储用户的签到历史记录，主要用于跨城市算法需要的用户源城市签到数据。

\begin{table}[hbt]
  \centering
  \caption{TripManner核心数据表设计}
  \label{tab:ch5_db}
  \begin{tabular*}{0.95\textwidth}{@{\extracolsep{\fill}}ccccc}
  \toprule
    数据表 & 核心字段 & 关联关系 & 索引 & 用途 \\
  \midrule
    \texttt{users} & id, username, password\_hash, preferences & --- & username唯一 & 用户信息与偏好 \\
    \texttt{pois} & id, name, category, city, lat, lng & --- & city索引 & POI信息库 \\
    \texttt{routes} & id, user\_id, algorithm\_used, route\_result & users(1:N) & user\_id索引 & 规划记录 \\
    \texttt{dialogs} & id, user\_id, title, messages & users(1:N) & user\_id索引 & 对话记录 \\
%    \texttt{mock\_routes} & id, city, algorithm\_type, route\_result & --- & city索引 & Mock预设数据 \\
    \texttt{user\_checkins} & id, user\_id, poi\_id, city, checkin\_time & users,pois & user\_id,city & 签到历史 \\
  \bottomrule
  \end{tabular*}
\end{table}

%Mock数据的设计覆盖了系统支持的全部数据集城市。单城市数据（EKD-Trip算法）为Glasgow、Osaka、Toronto和Tokyo每个城市各准备3条预设路线，覆盖不同的旅行模式类型（接近、远离、U型、不规则），共计12条。跨城市数据（CrossCityLLMCPR算法）为New York→Los Angeles、New York→San Francisco和Los Angeles→San Francisco三个城市对各准备2条预设路线，覆盖不同的偏好迁移场景，共计6条。全部18条Mock路线数据能够支撑系统演示的全部场景需求。

\section{系统实现}
\label{sec:ch5_implementation}

本系统基于上述架构和详细设计，采用前后端分离的方式进行开发实现。开发环境配置如下：操作系统为Windows 11，前端使用Node.js v18与npm进行包管理和构建，后端使用Python 3.11与pip进行依赖管理，数据库采用MySQL 8.0，开发工具为Visual Studio Code。

系统主页采用三栏式布局，如图\ref{fig:ch5_impl_main}所示。左侧为对话历史列表，中间为对话交互区域，右侧为Leaflet地图展示区域。对话交互区域支持流式文本渲染，路线规划结果以打字机效果逐字出现；地图区域配合流式输出进行动态路线绘制。在对话区域下方，意图可视化面板根据算法类型条件渲染不同的图表内容。底部输入框支持自然语言输入，旁边显示当前使用的算法标识。

\begin{figure}[hbt]
  \centering
  % TODO: 插入系统主页截图
  % 内容：三栏布局——左侧对话历史列表、中间对话区（含流式文本）、右侧Leaflet地图（含路线）、下方意图可视化面板、底部输入框
  \fbox{\parbox{0.85\textwidth}{\centering\vspace{4cm}系统主页界面截图占位\vspace{4cm}}}
  \caption{TripManner系统主页界面}
  \label{fig:ch5_impl_main}
\end{figure}

图\ref{fig:ch5_impl_ekd}展示了单城市路线规划场景下的系统界面。当用户查询东京一日游时，系统通过路由决策引擎识别为单城市场景并调用EKD-Trip算法。对话区域以流式方式逐步输出各景点的介绍文本，地图上同步动态添加POI标记和路线连线。路线输出完毕后，意图可视化面板以淡入效果显示，包含旅行模式图标（此例为接近模式）、距离变化曲线折线图以及模式说明文字。

\begin{figure}[hbt]
  \centering
  % TODO: 插入单城市路线+意图可视化截图
  % 内容：东京路线在地图上的展示 + 对话区的路线文本 + 意图面板（Travel Mode图标、距离曲线、说明文字）
  \fbox{\parbox{0.85\textwidth}{\centering\vspace{4cm}单城市路线及意图可视化截图占位\vspace{4cm}}}
  \caption{单城市路线规划与旅行模式可视化（EKD-Trip）}
  \label{fig:ch5_impl_ekd}
\end{figure}

图\ref{fig:ch5_impl_cross}展示了跨城市路线规划场景下的系统界面。当用户在纽约拥有签到历史并查询洛杉矶旅行时，系统调用CrossCity算法。意图可视化面板展示四维偏好因子雷达图、偏好迁移权重箭头图、城市群体偏好饼图以及个人与群体偏好融合比例进度条。用户可以直观地看到系统如何将其在纽约的偏好迁移到洛杉矶，以及最终路线中个人偏好和洛杉矶群体偏好各占多大比重。

\begin{figure}[hbt]
  \centering
  % TODO: 插入跨城市偏好迁移截图
  % 内容：洛杉矶路线在地图上 + 意图面板（雷达图、迁移权重图、群体偏好饼图、融合比例条）
  \fbox{\parbox{0.85\textwidth}{\centering\vspace{4cm}跨城市偏好迁移可视化截图占位\vspace{4cm}}}
  \caption{跨城市路线规划与偏好迁移可视化（CrossCityLLMCPR）}
  \label{fig:ch5_impl_cross}
\end{figure}

图\ref{fig:ch5_impl_llm}展示了通用城市场景下的流式输出效果。当用户查询北京旅行时，系统识别该城市不在算法数据集覆盖范围内，调用DeepSeek大语言模型Agent实时生成路线。由于是真实的大语言模型API调用，流式输出的打字机效果更加自然明显。地图上同步进行动态路线绘制，但意图可视化面板不予显示，因为大语言模型不产出结构化的意图推断数据。

\begin{figure}[hbt]
  \centering
  % TODO: 插入LLM Agent流式输出截图
  % 内容：对话区的打字机效果文本 + 地图上的北京路线 + 无意图面板
  \fbox{\parbox{0.85\textwidth}{\centering\vspace{4cm}LLM Agent流式路线规划截图占位\vspace{4cm}}}
  \caption{通用城市路线规划的流式输出（DeepSeek Agent）}
  \label{fig:ch5_impl_llm}
\end{figure}

%图\ref{fig:ch5_impl_admin}展示了系统管理页面，管理员可通过该页面切换Mock/Real运行模式，查看各算法支持的数据集城市列表以及系统当前的运行状态。
%
%\begin{figure}[hbt]
%  \centering
%  % TODO: 插入系统管理页截图
%  % 内容：Mock/Real切换开关 + 数据集城市列表 + 系统状态信息
%  \fbox{\parbox{0.85\textwidth}{\centering\vspace{3cm}系统管理页截图占位\vspace{3cm}}}
%  \caption{系统管理页——Mock/Real模式切换}
%  \label{fig:ch5_impl_admin}
%\end{figure}

\section{系统测试}
\label{sec:ch5_testing}

在系统开发和部署过程中，测试是确保稳定性和功能完备性的必要环节。本测试方案涵盖功能测试和性能测试两方面，验证系统能否按设计要求完成各项任务。

\subsection{测试环境配置}
\label{sec:ch5_test_env}

本系统测试运行环境配置如下：

\begin{itemize}
  \item 操作系统：Windows 11 Enterprise
  \item 数据库系统：MySQL 8.0
  \item 开发语言：Python 3.11（后端），JavaScript ES6+（前端）
  \item 前端框架：Vue.js 3，Vite，Element Plus，Leaflet，ECharts
  \item 后端框架：FastAPI，SQLAlchemy
  \item 测试工具：PyTest（单元测试），Postman（接口测试），Chrome DevTools（前端调试）
\end{itemize}

\subsection{功能测试}
\label{sec:ch5_func_test}

功能测试围绕系统的五个核心演示场景设计测试用例，覆盖三种算法路径、引导式对话等关键功能。测试结果如表\ref{tab:ch5_functest}所示。

\begin{table}[hbt]
  \centering
  \caption{功能测试结果}
  \label{tab:ch5_functest}
  \begin{tabular*}{0.95\textwidth}{@{\extracolsep{\fill}}cccc}
  \toprule
    测试场景 & 测试点 & 测试预期 & 测试结果 \\
  \midrule
    单城市-Tokyo & 引导式对话 & 缺失字段时逐步追问 & 通过 \\
    单城市-Tokyo & EKD-Trip路由 & 正确调用EKD-Trip算法 & 通过 \\
    单城市-Tokyo & 旅行模式可视化 & 显示接近模式与距离曲线 & 通过 \\
    单城市-Glasgow & 流式输出 & 文本逐字渲染、地图动态绘制 & 通过 \\
    单城市-Glasgow & U型模式展示 & 显示U型模式图标与曲线 & 通过 \\
    跨城市-NY→LA & 跨城市路由 & 正确调用CrossCity & 通过 \\
    跨城市-NY→LA & 偏好雷达图 & 展示4维偏好因子分布 & 通过 \\
    跨城市-NY→LA & 迁移权重展示 & 展示偏好迁移系数 & 通过 \\
    通用城市-北京 & DeepSeek Agent & 实时LLM流式生成路线 & 通过 \\
    通用城市-北京 & 无意图面板 & 不显示意图可视化 & 通过 \\
    用户管理 & 注册登录 & 认证正常运行 & 通过 \\
  \bottomrule
  \end{tabular*}
\end{table}

功能测试结果表明，系统各功能模块均按设计要求正常运行。路由决策引擎根据目的城市正确选择了算法路径，引导式对话有效补全了不完整的查询信息，流式输出和动态路线渲染运行流畅，意图可视化面板根据算法类型正确渲染了对应图表。

\subsection{性能测试}
\label{sec:ch5_perf_test}

性能测试主要验证系统在不同运行模式下的响应时间是否满足非功能性需求中的指标要求。测试结果如表\ref{tab:ch5_perftest}所示。

\begin{table}[hbt]
  \centering
  \caption{性能测试结果}
  \label{tab:ch5_perftest}
  \begin{tabular*}{0.92\textwidth}{@{\extracolsep{\fill}}cccc}
  \toprule
    测试项目 & 指标要求 & 实测结果 & 是否达标 \\
  \midrule
    页面首次加载 & $<$ 3s & 1.8s & 达标 \\
    LLM Agent首个SSE事件 & $<$ 2s & 1.2s & 达标 \\
    流式文本逐Token延迟 & $<$ 100ms & 42ms & 达标 \\
    POI标记动态添加 & 300-800ms & 500ms & 达标 \\
    算法路线完整输出 & $<$ 5s & 2.1s & 达标 \\
    LLM路线完整输出 & $<$ 15s & 8.6s & 达标 \\
  \bottomrule
  \end{tabular*}
\end{table}

性能测试结果表明，系统各项响应时间指标均满足非功能性需求中的要求。算法路线规划响应迅速，完整路线输出仅需2.1秒；DeepSeek Agent模式下由于涉及外部API调用，首个SSE事件延迟略高但仍在2秒以内，完整路线输出耗时8.6秒，用户在流式输出过程中即可逐步看到路线生成结果，无需等待全部完成后才能查看。

\section{本章小结}
\label{sec:ch5_summary}

本章设计并实现了TripManner智能旅行规划系统，将第三章和第四章提出的路线推荐算法应用于实际旅行规划场景。系统采用Vue 3与FastAPI的前后端分离架构，集成EKD-Trip单城市算法、CrossCityLLMCPR跨城市算法和DeepSeek大语言模型Agent三条算法路径，由路由决策引擎根据目的城市自动选择路线规划方式。

系统在两个方面有所创新：一方面，意图可视化模块将算法推理过程中的用户偏好意图（旅行模式、多维偏好因子、偏好迁移权重、城市群体偏好等）以图表形式呈现，使用户不仅能获得推荐路线，还能理解推荐背后的逻辑，从而增强对系统的信任。另一方面，引导式对话补全机制和SSE流式输出带来了自然流畅的交互体验——用户无需一次性提供完整查询信息，系统会逐步引导补全，路线结果以打字机效果呈现并在地图上动态绘制。

功能测试和性能测试结果表明，系统各功能模块运行正常，响应时间指标满足设计要求，能够支撑流畅的用户交互体验。
